
%% -------------------------------------------------------------------
\subsection{\vrev{방법론적 기여}}
\label{subsec:contrib-methodological}
%% -------------------------------------------------------------------

\fdel{방법론적 관점에서 본 연구는 다음의 세 가지 기여를 달성하였다.}

\fdel{첫째, \textbf{P85.8 기준 분할 회귀(split regression)를 파레토 분포 분석에 적용}하였다.
C3 전체 표본에서 2차교호작용 회귀의 $R^2 = 0.858$에 그치는 잔여 비설명 분산을,
P85.8 기준으로 Core($\leq$ P85.8)와 Tail($>$ P85.8)로 분할함으로써
Core $R^2 = 0.999$(0.9990), Tail $R^2 = 0.710$(0.7100)으로 Core 영역의
결정론적 구조를 완전히 포착하였다.
이 분할 전략은 파레토 분포의 체제 전환(regime change)을 포착하는
범용적 분석 방법론으로, BCM을 넘어 보험, 금융, 재난관리 등
파레토 현상이 핵심인 분야에 확장 적용될 수 있다.}

\fdel{둘째, \textbf{비중첩 대역 시드(non-overlapping band seeding) 체계}를 설계하였다.
마스터 시드 $s = 42$로부터 변수별 시드를 $10{,}000$ 간격의 비중첩 대역($O_{\mathrm{raw}}$: $[0\text{--}9999]$, $k_O$: $[20000\text{--}29999]$,
$U^{(G)}$: $[30000\text{--}39999]$,
$U^{(P)}$: $[40000\text{--}49999]$)으로 파생하여,
시뮬레이션의 완전한 재현성을 보장하면서도 변수 간 통계적 독립성을
유지한다(모든 변수 쌍 $|r| < 0.02$).}

\jrev{넷째}, \textbf{6단계 다층 검증(six-stage multi-layer validation) 프레임워크}를 구축하였다.
기반 검증(L0, Golden Compare), 시뮬레이션 검증, 회귀 검증, 강건성 검증, 이중채널 검증, 전문가 평가의
6개 단계(L0--L5)가 계층적 의존 구조를 형성하여 모델의 타당성을 다각적으로 보장한다.
정량적 시뮬레이션과 정성적 전문가 평가를 삼각검증(triangulation)함으로써
단일 방법론의 편향을 최소화하였으며, 이 프레임워크는
수리 모델 기반 연구의 검증 설계 표준으로 활용될 수 있다.


